% =============================================================================
% Chapter 14: Audit Findings and Resolutions
% =============================================================================
\chapter{Audit Findings and Resolutions}
\label{ch:audit-findings}

This chapter documents audit findings from the formal review and their 
resolutions. All P0 (critical) and P1 (high priority) findings are 
addressed before production release.

% =============================================================================
\section{Audit Review Summary}
\label{sec:audit-summary}
% =============================================================================

\begin{tcolorbox}[colback=yellow!10,colframe=yellow!50!black,title=Audit Status: Conditionally Approved]
\textbf{Review Date:} 20 April 2026\\
\textbf{Overall Verdict:} Conditionally Approved\\
\textbf{Condition:} P0 items must be resolved before GA cutover

\vspace{0.5em}
\begin{tabular}{lccl}
\toprule
\textbf{Checkpoint} & \textbf{Status} & \textbf{Condition} \\
\midrule
H1 - Source Completeness & \textcolor{green}{\textbf{PASS}} & None \\
H2 - Extraction Accuracy & \textcolor{green}{\textbf{PASS}} & None \\
H3 - Requirement Derivation & \textcolor{orange}{\textbf{CONDITIONAL}} & Pending G-001 resolution \\
H4 - Traceability Links & \textcolor{green}{\textbf{PASS}} & None \\
H5 - Schema Alignment & \textcolor{orange}{\textbf{CONDITIONAL}} & Pending audit trail schema \\
H6 - Gap Analysis & \textcolor{orange}{\textbf{CONDITIONAL}} & Pending G-003 cross-spec \\
\bottomrule
\end{tabular}
\end{tcolorbox}

% =============================================================================
\section{P0 Findings: Critical (Resolution Required Before v1.0)}
\label{sec:p0-findings}
% =============================================================================

% -----------------------------------------------------------------------------
\subsection{P0-1: State Count Reconciliation (27 vs 28)}
\label{sec:p0-state-count}
% -----------------------------------------------------------------------------

\begin{tcolorbox}[colback=red!5,colframe=red!70,title=\faIcon{exclamation-triangle} Finding P0-1: State Count Inconsistency]
\textbf{Location:} HITL Spec p.9 vs Implementation Spec Table 5.1\\
\textbf{Issue:} HITL spec references ``27 states + 3 error-handling'' while 
implementation spec Table 5.1 enumerates 28 states\\
\textbf{Severity:} \textbf{P0 - Critical}\\
\textbf{Owner:} Spec Owner\\
\textbf{Timeline:} 1 week
\end{tcolorbox}

\textbf{Resolution:}

\begin{tcolorbox}[colback=green!5,colframe=green!70,title=\faIcon{check-circle} Resolution Applied]
The \textbf{implementation spec's 28-state enumeration is authoritative}.

The discrepancy arose from counting methodology:
\begin{itemize}
    \item \textbf{HITL Spec (27+3):} Counted 27 ``primary'' states plus 3 error-handling 
          states as a separate category = 30 total states
    \item \textbf{Implementation Spec (28):} Counted 28 operational states; error-handling 
          states (AWAITING\_CLARIFICATION, ESCALATED\_TO\_FINANCE, MANUAL\_OVERRIDE) 
          are subsumed within the 28
\end{itemize}

\textbf{Authoritative Source:} \texttt{structured/workflows/tb-review-glo.json}\\
\textbf{Correct Count:} \textbf{28 states} as enumerated in Table 5.1

\vspace{0.5em}
\textbf{Corrective Action:} This specification is updated to reference 28 states 
throughout. Cross-references corrected in:
\begin{itemize}
    \item Chapter~\ref{ch:source-documents}: Updated BPMN extraction to state ``28 states''
    \item Chapter~\ref{ch:specification-mapping}: Updated process step count
    \item YAML manifest: Updated \texttt{statistics.bpmn\_states} to 28
\end{itemize}
\end{tcolorbox}

\textbf{State Enumeration (Authoritative):}

\begin{longtable}{clll}
\caption{Authoritative 28-State Enumeration} \\
\toprule
\textbf{\#} & \textbf{State ID} & \textbf{Type} & \textbf{Timing} \\
\midrule
\endfirsthead
\toprule
\textbf{\#} & \textbf{State ID} & \textbf{Type} & \textbf{Timing} \\
\midrule
\endhead
1 & ROLL\_FORWARD & Task & M-3 \\
2 & DOWNLOAD\_PSGL & Task & M-3 \\
3 & UPDATE\_ACCOUNTS\_MASTER & Task & M-3 \\
4 & EXTRACT\_TB & Task & M-3 to M+5 \\
5 & UPDATE\_COMPARATIVES & Task & M-3 to M+5 \\
6 & CALC\_VARIANCE & Task & M-2 to M+5 \\
7 & FILTER\_VARIANCES & Task & M-2 to M+5 \\
8 & UPDATE\_COMMENTARIES & Task & M-2 to M+5 \\
9 & CHECK\_JOURNAL\_POSTED & Gateway & M-2 to M+5 \\
10 & SEND\_JOURNAL\_REQUEST & Task & M-2 to M+5 \\
11 & UPDATE\_RISK\_SUMMARY & Task & M-2 to M+5 \\
12 & PREPARE\_SUMMARY & Task & M-2 to M+5 \\
13 & INVESTIGATE\_VARIANCE & Task & M-2 to M+5 \\
14 & CHECK\_MATERIAL & Gateway & M-2 to M+5 \\
15 & INTERNAL\_REVIEW & Task & M+5 \\
16 & SEND\_FORTM\_SUMMARY & Task & M+5 \\
17 & CATEGORIZE\_UNEXPLAINED & Task & M-2 to M+5 \\
18 & CHECK\_TEAM\_ENTRIES & Task & M-2 to M+5 \\
19 & CONFIRM\_ENTRIES & Task & M-2 to M+5 \\
20 & SEEK\_COMMENTARIES & Task & M-2 to M+5 \\
21 & RECEIVE\_COMMENTARIES & Task & M-2 to M+5 \\
22 & PERFORM\_RISK\_ANALYSIS & Task & M-2 to M+5 \\
23 & UPDATE\_VARIANCE\_COMMENTARIES & Task & M-2 to M+5 \\
24 & TRACK\_ENTRIES & Task & M-2 to M+5 \\
25 & AWAITING\_CLARIFICATION & Error & As needed \\
26 & ESCALATED\_TO\_FINANCE & Error & As needed \\
27 & MANUAL\_OVERRIDE & Error & As needed \\
28 & COMPLETED & End & M+5 \\
\bottomrule
\end{longtable}

% -----------------------------------------------------------------------------
\subsection{P0-2: vLLM Fallback Specification}
\label{sec:p0-vllm-fallback}
% -----------------------------------------------------------------------------

\begin{tcolorbox}[colback=red!5,colframe=red!70,title=\faIcon{exclamation-triangle} Finding P0-2: vLLM Fallback Not Specified]
\textbf{Location:} tb-review-spec, p.13, \S3.1\\
\textbf{Issue:} ``vllm-only routing policy'' with no fallback if vLLM cluster degrades\\
\textbf{Impact:} Commentary generation (Stage 3) blocks on LLM availability\\
\textbf{Severity:} \textbf{P0 - Critical}\\
\textbf{Owner:} AI Architect\\
\textbf{Timeline:} 2 weeks
\end{tcolorbox}

\textbf{Resolution:}

\begin{tcolorbox}[colback=green!5,colframe=green!70,title=\faIcon{check-circle} Resolution: Graceful Degradation Specification]
\textbf{Primary:} vLLM cluster via AI Core (default)\\
\textbf{Fallback:} Template-based commentary generation\\
\textbf{Trigger Conditions:} (any of)
\begin{itemize}
    \item vLLM latency P95 $> 5000$ms (5-minute rolling window)
    \item vLLM error rate $> 10\%$ (5-minute rolling window)
    \item vLLM cluster health check fails 3 consecutive probes
\end{itemize}
\end{tcolorbox}

\textbf{Fallback Implementation Specification:}

\begin{lstlisting}[language=Python,caption={vLLM Fallback Logic}]
class CommentaryService:
    VLLM_LATENCY_THRESHOLD_MS = 5000
    VLLM_ERROR_RATE_THRESHOLD = 0.10
    HEALTH_CHECK_FAILURES_MAX = 3
    
    def generate_commentary(self, variance_record: VarianceRecord) -> Commentary:
        if self._should_use_fallback():
            return self._generate_template_commentary(variance_record)
        try:
            return self._generate_vllm_commentary(variance_record)
        except VLLMTimeoutError:
            self._record_latency_breach()
            return self._generate_template_commentary(variance_record)
    
    def _generate_template_commentary(self, vr: VarianceRecord) -> Commentary:
        """Template-based fallback with 80% of LLM quality."""
        template = self._select_template(vr.variance_type, vr.account_category)
        return Commentary(
            text=template.format(
                account=vr.account_name,
                variance_amount=format_currency(vr.variance_amount),
                variance_pct=format_percent(vr.variance_percentage),
                direction="increased" if vr.variance_amount > 0 else "decreased",
                prior_period=vr.prior_period,
            ),
            confidence=0.75,  # Template confidence capped at 75%
            source="template_fallback",
            llm_available=False,
        )
\end{lstlisting}

\textbf{Template Library Requirements:}

\begin{tabular}{lll}
\toprule
\textbf{Template ID} & \textbf{Variance Type} & \textbf{Account Category} \\
\midrule
TPL-001 & Spike ($>25\%$) & Revenue \\
TPL-002 & Spike ($>25\%$) & Cost of Sales \\
TPL-003 & Drop ($<-25\%$) & Revenue \\
TPL-004 & Drop ($<-25\%$) & Cost of Sales \\
TPL-005 & New Account & Any \\
TPL-006 & Material Balance & Any \\
TPL-007 & General Variance & Any (catch-all) \\
\bottomrule
\end{tabular}

\textbf{Monitoring Requirements:}

\begin{enumerate}
    \item Dashboard SHALL display ``LLM Status'' indicator (green/amber/red)
    \item Fallback activations SHALL be logged with timestamp and trigger reason
    \item Weekly report SHALL include fallback activation count and duration
    \item SLA: Fallback mode SHALL not exceed 4 hours continuous without escalation
\end{enumerate}

% =============================================================================
\section{P1 Findings: High Priority (Resolution Required Within 4 Weeks)}
\label{sec:p1-findings}
% =============================================================================

% -----------------------------------------------------------------------------
\subsection{P1-1: Audit Trail Schema for EUC Retirement}
\label{sec:p1-audit-trail}
% -----------------------------------------------------------------------------

\begin{tcolorbox}[colback=orange!5,colframe=orange!70,title=\faIcon{exclamation-circle} Finding P1-1: Audit Trail Schema Not Defined]
\textbf{Location:} tb-review-spec, Section 1.8 (p.6-7)\\
\textbf{Issue:} Dashboard audit trail defined as ``authoritative evidence'' but no 
field-level schema specified\\
\textbf{Impact:} Cannot verify EUC retirement evidence standard\\
\textbf{Severity:} \textbf{P1 - High}\\
\textbf{Owner:} Controls Owner\\
\textbf{Timeline:} 4 weeks
\end{tcolorbox}

\textbf{Resolution:}

\begin{tcolorbox}[colback=green!5,colframe=green!70,title=\faIcon{check-circle} Schema Definition: audit-trail-entry.schema.json]
\begin{lstlisting}[language=jsonx,basicstyle=\tiny\ttfamily]
{
  "$schema": "https://json-schema.org/draft/2020-12/schema",
  "$id": "audit-trail-entry.schema.json",
  "title": "Audit Trail Entry",
  "type": "object",
  "required": ["audit_id", "timestamp", "persona_id", "action", "evidence_hash"],
  "properties": {
    "audit_id": {
      "type": "string",
      "format": "uuid",
      "description": "Unique identifier for this audit entry"
    },
    "timestamp": {
      "type": "string",
      "format": "date-time",
      "description": "ISO 8601 timestamp of the action"
    },
    "persona_id": {
      "type": "string",
      "description": "Employee ID from SAP IAS token (sub claim)"
    },
    "action": {
      "type": "string",
      "enum": ["view", "claim", "decide", "escalate", "override", "bulk_assign", 
               "comment_edit", "comment_approve", "comment_reject", "export"]
    },
    "resource_type": {
      "type": "string",
      "enum": ["variance_record", "commentary", "workflow_state", "entity_params"]
    },
    "resource_id": {
      "type": "string",
      "description": "ID of the affected resource"
    },
    "before_state": {
      "type": "object",
      "description": "Snapshot of resource state before action (null for create)"
    },
    "after_state": {
      "type": "object",
      "description": "Snapshot of resource state after action"
    },
    "evidence_hash": {
      "type": "string",
      "pattern": "^sha256:[a-f0-9]{64}$",
      "description": "SHA-256 hash of (timestamp + persona_id + action + resource_id + after_state)"
    },
    "bulk_operation_id": {
      "type": "string",
      "format": "uuid",
      "description": "Groups entries from same bulk operation (null for single ops)"
    },
    "ip_address": {
      "type": "string",
      "format": "ipv4",
      "description": "Client IP address"
    },
    "user_agent": {
      "type": "string",
      "description": "Browser/client user agent string"
    }
  }
}
\end{lstlisting}
\end{tcolorbox}

\textbf{Divergence Detection Specification:}

During the 12-month EUC deprecation window, divergence between Excel EUC and 
dashboard may occur. The following automated detection is specified:

\begin{lstlisting}[language=SQL,caption={Daily Divergence Detection Query}]
-- Run daily at M+6 for prior month data
WITH euc_data AS (
    SELECT legal_entity, account, closing_balance 
    FROM imported_euc_tb WHERE period = :period
),
dashboard_data AS (
    SELECT legal_entity, account, closing_balance 
    FROM variance_records WHERE period = :period
)
SELECT 
    COALESCE(e.legal_entity, d.legal_entity) AS le,
    COALESCE(e.account, d.account) AS account,
    e.closing_balance AS euc_balance,
    d.closing_balance AS dashboard_balance,
    ABS(e.closing_balance - d.closing_balance) AS divergence,
    CASE 
        WHEN e.closing_balance IS NULL THEN 'dashboard_only'
        WHEN d.closing_balance IS NULL THEN 'euc_only'
        ELSE 'value_mismatch'
    END AS divergence_type
FROM euc_data e
FULL OUTER JOIN dashboard_data d 
    ON e.legal_entity = d.legal_entity AND e.account = d.account
WHERE e.closing_balance != d.closing_balance 
   OR e.closing_balance IS NULL OR d.closing_balance IS NULL;
\end{lstlisting}

\textbf{Resolution Rule:} Dashboard is authoritative. Divergence records are 
logged with \texttt{action="divergence\_detected"} and require manual review 
by Controls Owner within 5 business days.

% -----------------------------------------------------------------------------
\subsection{P1-2: Statistical Normality Test for 3$\sigma$ Threshold}
\label{sec:p1-normality}
% -----------------------------------------------------------------------------

\begin{tcolorbox}[colback=orange!5,colframe=orange!70,title=\faIcon{exclamation-circle} Finding P1-2: Statistical Assumption Not Validated]
\textbf{Location:} tb-review-spec, p.14, \S3.3\\
\textbf{Issue:} 3$\sigma$ threshold assumes normally-distributed variance; 
financial data often exhibits fat tails (leptokurtosis)\\
\textbf{Severity:} \textbf{P1 - High}\\
\textbf{Owner:} Data Science\\
\textbf{Timeline:} 4 weeks
\end{tcolorbox}

\textbf{Resolution:}

\begin{tcolorbox}[colback=green!5,colframe=green!70,title=\faIcon{check-circle} Resolution: Normality Testing with MAD Fallback]
\textbf{Acceptance Criteria Addition (AC-ES01-2a):}

Before applying Z-score $> 3.0$ threshold, the system SHALL:
\begin{enumerate}
    \item Run Shapiro-Wilk normality test per account category
    \item If $p$-value $> 0.05$: Use Z-score (normal distribution confirmed)
    \item If $p$-value $\leq 0.05$: Use MAD (Median Absolute Deviation) with 
          threshold $= 3.5 \times$ MAD (robust to fat tails)
\end{enumerate}
\end{tcolorbox}

\textbf{Implementation:}

\begin{lstlisting}[language=Python,caption={Normality Test with MAD Fallback}]
from scipy import stats
import numpy as np

class AnomalyDetector:
    ZSCORE_THRESHOLD = 3.0
    MAD_THRESHOLD = 3.5
    SHAPIRO_P_THRESHOLD = 0.05
    
    def detect_anomalies(self, values: np.ndarray) -> np.ndarray:
        """Returns boolean mask of anomalies."""
        # Test for normality
        if len(values) >= 20:  # Shapiro-Wilk requires n >= 20
            _, p_value = stats.shapiro(values)
            use_robust = p_value <= self.SHAPIRO_P_THRESHOLD
        else:
            use_robust = True  # Default to robust for small samples
        
        if use_robust:
            return self._mad_detection(values)
        else:
            return self._zscore_detection(values)
    
    def _zscore_detection(self, values: np.ndarray) -> np.ndarray:
        z_scores = np.abs(stats.zscore(values))
        return z_scores > self.ZSCORE_THRESHOLD
    
    def _mad_detection(self, values: np.ndarray) -> np.ndarray:
        median = np.median(values)
        mad = np.median(np.abs(values - median))
        # Consistency constant for normal distribution approximation
        mad_adjusted = mad * 1.4826
        modified_z = np.abs(values - median) / mad_adjusted
        return modified_z > self.MAD_THRESHOLD
\end{lstlisting}

\textbf{Documentation Note:} Add to Chapter~\ref{ch:business-requirements}, 
AC-ES01-2: ``Anomaly detection SHALL flag items with Z-score $> 3.0$ for 
normally-distributed data (Shapiro-Wilk $p > 0.05$) \textbf{OR} MAD-score 
$> 3.5$ for non-normal distributions.''

% =============================================================================
\section{Cross-Cutting Issue Resolutions}
\label{sec:cross-cutting}
% =============================================================================

% -----------------------------------------------------------------------------
\subsection{C-1: Terminology Reconciliation}
\label{sec:c1-terminology}
% -----------------------------------------------------------------------------

\begin{tabular}{lll}
\toprule
\textbf{Term} & \textbf{Authoritative Value} & \textbf{Corrective Action} \\
\midrule
State count & \textbf{28 states} & Update HITL spec to 28 \\
EUC-002 name & ``TB Variance Analysis file integrity'' & Standardize across docs \\
\bottomrule
\end{tabular}

% -----------------------------------------------------------------------------
\subsection{C-2: Materiality Threshold Standardization}
\label{sec:c2-materiality}
% -----------------------------------------------------------------------------

\begin{tcolorbox}[colback=blue!5,colframe=blue!70,title=Resolution: Entity-Params Configuration]
\textbf{Issue:} Hardcoded thresholds (BS $> \$100$M, PL $> \$3$M) in DP-001 
conflict with ``country-specific'' in TB-KPI-001.

\textbf{Resolution:} Move thresholds to \texttt{entity-params/<le-code>.yaml}:
\begin{lstlisting}[language=yaml]
# entity-params/HKG.yaml
legal_entity_code: "HKG"
materiality_thresholds:
  balance_sheet_absolute: 100000000  # $100M USD
  profit_loss_absolute: 3000000      # $3M USD
  variance_percentage: 0.10          # 10%
\end{lstlisting}

\textbf{Corrective Action:} Remove hardcoded thresholds from DP-001 description; 
reference ``entity-specific materiality from \texttt{entity-params}'' instead.
\end{tcolorbox}

% =============================================================================
\section{Additional Auditor Findings}
\label{sec:additional-findings}
% =============================================================================

% -----------------------------------------------------------------------------
\subsection{A-2: 100\% Traceability Claim Qualification}
\label{sec:a2-traceability-exceptions}
% -----------------------------------------------------------------------------

\begin{tcolorbox}[colback=orange!5,colframe=orange!70,title=\faIcon{exclamation-circle} Finding A-2: Exceptions to 100\% Traceability]
\textbf{Issue:} The spec claims ``100\% BIDIRECTIONAL TRACEABILITY ACHIEVED'' 
but exceptions exist and should be explicitly disclosed.
\end{tcolorbox}

\textbf{Resolution: Exceptions Table}

\begin{longtable}{p{3cm} p{3.5cm} p{6cm}}
\caption{Explicit Exceptions to 100\% Traceability} \\
\toprule
\textbf{Layer} & \textbf{Item} & \textbf{Justification} \\
\midrule
\endfirsthead
\toprule
\textbf{Layer} & \textbf{Item} & \textbf{Justification} \\
\midrule
\endhead
Spec $\rightarrow$ Schema & TB-STEP-006 (Journal Posting Investigation) & Workflow orchestration step; no persistent data schema required \\
Spec $\rightarrow$ Schema & TB-STEP-007 (Track Entries \& Risk) & Workflow orchestration step; state tracked in workflow engine, not schema \\
Schema $\rightarrow$ Spec & \texttt{audit\_id} field & Infrastructure field for tamper detection; not derived from business requirement \\
Schema $\rightarrow$ Spec & \texttt{version} field & Infrastructure field for concurrency control; not derived from business requirement \\
Schema $\rightarrow$ Spec & \texttt{processed\_at} field & Infrastructure field for pipeline telemetry; not derived from business requirement \\
\bottomrule
\end{longtable}

\textbf{Updated Claim:}
\begin{quote}
\textit{``100\% BIDIRECTIONAL TRACEABILITY ACHIEVED for business requirements. 
Five infrastructure/orchestration items are documented exceptions (see Table 14.X).''}
\end{quote}

% -----------------------------------------------------------------------------
\subsection{A-4: BSS Risk Assessment Vocabulary Alignment}
\label{sec:a4-bss-vocabulary}
% -----------------------------------------------------------------------------

\begin{tcolorbox}[colback=orange!5,colframe=orange!70,title=\faIcon{exclamation-circle} Finding A-4: BSS Cross-Spec Vocabulary]
\textbf{Issue:} BSS Risk Assessment is out of scope but risk vocabulary 
(\texttt{substantiated\_at\_risk}, \texttt{unsubstantiated\_unowned}, 
\texttt{unsubstantiated\_owned}) appears in DP-005.\\
\textbf{Question:} How will cross-spec vocabulary alignment be maintained?
\end{tcolorbox}

\textbf{Resolution: Shared Vocabulary Registry}

\begin{tcolorbox}[colback=green!5,colframe=green!70,title=\faIcon{check-circle} Resolution: Cross-Spec Enum Registry]
A shared vocabulary registry SHALL be maintained at:
\begin{verbatim}
docs/schema/common/enums.yaml
\end{verbatim}

Risk status values are defined once and imported by both TB and BSS specs:

\begin{lstlisting}[language=yaml,caption={Shared Enum Registry (enums.yaml)}]
# docs/schema/common/enums.yaml
risk_status:
  $id: "common/risk-status"
  description: "Shared risk classification vocabulary (TB + BSS)"
  type: string
  enum:
    - substantiated           # Evidence supports balance
    - substantiated_at_risk   # Evidence exists but control weakness identified
    - unsubstantiated_unowned # No evidence; no owner assigned
    - unsubstantiated_owned   # No evidence; owner assigned for remediation
  ownership:
    maintainer: "Controls Architecture Team"
    consumers: ["tb-review-spec", "bss-risk-assessment-spec"]
    change_process: "RFC required; both spec owners must approve"
\end{lstlisting}

\textbf{Governance:}
\begin{enumerate}
    \item Any change to shared enums requires RFC (Request for Comment)
    \item Both TB and BSS spec owners must approve vocabulary changes
    \item CI validation ensures specs reference registry (not local definitions)
\end{enumerate}
\end{tcolorbox}

% -----------------------------------------------------------------------------
\subsection{I-3: Bulk Operations Error Handling}
\label{sec:i3-bulk-operations}
% -----------------------------------------------------------------------------

\begin{tcolorbox}[colback=orange!5,colframe=orange!70,title=\faIcon{exclamation-circle} Finding I-3: Bulk Operations Gaps]
\textbf{Missing specifications:}
\begin{itemize}
    \item Partial success handling (e.g., 10 of 50 items fail)
    \item Race conditions (two reviewers claiming same item via bulk)
    \item Audit of bulk operations (single entry or per-item?)
\end{itemize}
\end{tcolorbox}

\textbf{Resolution: Bulk Operations Specification}

\begin{tcolorbox}[colback=green!5,colframe=green!70,title=\faIcon{check-circle} Resolution: Bulk Operation Semantics]

\textbf{1. Atomic Per-Item with Rollback:}
\begin{itemize}
    \item Each item in a bulk operation SHALL be processed atomically
    \item If \textbf{any} item fails, the \textbf{entire bulk operation} SHALL be rolled back
    \item No partial success state is allowed
\end{itemize}

\textbf{2. Race Condition Prevention:}
\begin{lstlisting}[language=SQL,caption={Optimistic Locking for Bulk Claim}]
-- Bulk claim with optimistic locking
BEGIN TRANSACTION;
  -- Attempt to claim all items atomically
  UPDATE variance_records 
  SET claimant_id = :persona_id,
      claimed_at = NOW(),
      version = version + 1
  WHERE id IN (:item_ids)
    AND claimant_id IS NULL  -- Only unclaimed items
    AND version = :expected_version;  -- Optimistic lock
  
  -- Check all items were claimed
  IF ROW_COUNT() != :expected_count THEN
    ROLLBACK;  -- Race condition detected
    RETURN error: 'CONFLICT: Some items already claimed'
  END IF;
COMMIT;
\end{lstlisting}

\textbf{3. Audit of Bulk Operations:}
\begin{itemize}
    \item Each item SHALL have its own audit entry
    \item All entries SHALL share the same \texttt{bulk\_operation\_id}
    \item A summary audit entry SHALL be created with operation context
\end{itemize}

\begin{lstlisting}[language=jsonx,caption={Bulk Operation Audit Example}]
// Summary entry
{
  "audit_id": "uuid-summary",
  "action": "bulk_assign",
  "bulk_operation_id": "uuid-bulk-123",
  "bulk_operation_context": {
    "total_items": 50,
    "success_count": 50,
    "failure_count": 0,
    "operation_type": "bulk_assign"
  }
}

// Individual item entries (50 of these)
{
  "audit_id": "uuid-item-1",
  "action": "bulk_assign",
  "bulk_operation_id": "uuid-bulk-123",
  "resource_id": "variance-record-001",
  "change_summary": {
    "new_assignee": "1234567"
  }
}
\end{lstlisting}
\end{tcolorbox}

% -----------------------------------------------------------------------------
\subsection{I-4: Kill-Switch Automation Specification}
\label{sec:i4-kill-switch}
% -----------------------------------------------------------------------------

\begin{tcolorbox}[colback=orange!5,colframe=orange!70,title=\faIcon{exclamation-circle} Finding I-4: Kill-Switch Automation Missing]
\textbf{Table 11.4 defines six kill-switch conditions but doesn't specify:}
\begin{itemize}
    \item Who evaluates? (Automated monitor vs. manual check)
    \item Evaluation frequency (continuous? post-cycle?)
    \item Notification mechanism (email? PagerDuty? Jira?)
\end{itemize}
\end{tcolorbox}

\textbf{Resolution: Kill-Switch Automation Specification}

\begin{longtable}{p{2.5cm} p{2cm} p{2cm} p{3cm} p{3cm}}
\caption{Kill-Switch Condition Automation} \\
\toprule
\textbf{Condition} & \textbf{Evaluator} & \textbf{Frequency} & \textbf{Notification} & \textbf{Probe Implementation} \\
\midrule
\endfirsthead
\toprule
\textbf{Condition} & \textbf{Evaluator} & \textbf{Frequency} & \textbf{Notification} & \textbf{Probe Implementation} \\
\midrule
\endhead
Accept-without-edit $<$ 50\% & Automated & Post-cycle (M+6) & PagerDuty P2 + Slack & SQL query on decisions table \\
Material variance missed & Automated & Daily at 06:00 UTC & PagerDuty P1 + Email & Compare AI flags vs manual flags \\
Commentary factual error & Manual & Weekly review & Jira ticket & Random sampling (5\%) \\
Hallucination detected & Automated & Real-time & PagerDuty P1 + Kill & Confidence score $<$ 0.3 AND no source \\
vLLM availability $<$ 95\% & Automated & Continuous (1-min) & PagerDuty P2 & Health check endpoint \\
Audit trail integrity & Automated & Hourly & PagerDuty P1 + Kill & Evidence hash verification \\
\bottomrule
\end{longtable}

\textbf{Probe Implementation Examples:}

\begin{lstlisting}[language=Python,caption={Kill-Switch Probe Implementation}]
class KillSwitchMonitor:
    """Automated kill-switch condition evaluation."""
    
    def check_accept_rate(self, period: str) -> KillSwitchResult:
        """Post-cycle check: Accept-without-edit rate."""
        query = """
            SELECT 
                COUNT(*) FILTER (WHERE decision='accept' AND edits='{}') AS accept_no_edit,
                COUNT(*) AS total_reviewed
            FROM decisions WHERE period = %s
        """
        result = db.execute(query, [period])
        rate = result.accept_no_edit / result.total_reviewed
        if rate < 0.50:
            return KillSwitchResult(
                triggered=True,
                condition="accept_rate_below_50",
                value=rate,
                action="ALERT_P2",
                notification=["pagerduty:p2", "slack:#tb-review-alerts"]
            )
        return KillSwitchResult(triggered=False)
    
    def check_hallucination_realtime(self, commentary: Commentary) -> KillSwitchResult:
        """Real-time check: Hallucination detection."""
        if commentary.confidence < 0.3 and not commentary.source_references:
            return KillSwitchResult(
                triggered=True,
                condition="hallucination_detected",
                value=commentary.confidence,
                action="KILL_IMMEDIATELY",
                notification=["pagerduty:p1", "kill-switch:activate"]
            )
        return KillSwitchResult(triggered=False)
\end{lstlisting}

% =============================================================================
\section{Commentary Quality Metrics Clarification}
\label{sec:commentary-metrics}
% =============================================================================

\begin{tcolorbox}[colback=orange!5,colframe=orange!70,title=\faIcon{exclamation-circle} Finding P2-2: Metrics Ambiguity]
\textbf{Location:} tb-review-spec, p.15, \S3.8\\
\textbf{Issue:} ``accept-with-edit'' numerator/denominator not defined
\end{tcolorbox}

\textbf{Authoritative Definitions:}

\begin{align}
\text{accept\_without\_edit\_rate} &= \frac{|\{d : d.\text{decision} = \text{`accept'} \land d.\text{edits} = \emptyset\}|}{|\text{total\_reviewed}|} \geq 70\% \\[1em]
\text{accept\_with\_any\_edit\_rate} &= \frac{|\{d : d.\text{decision} = \text{`accept'}\}|}{|\text{total\_reviewed}|} \geq 95\%
\end{align}

\textbf{Interpretation:}
\begin{itemize}
    \item \textbf{accept\_without\_edit\_rate:} Commentaries accepted exactly as AI generated
    \item \textbf{accept\_with\_any\_edit\_rate:} All accepted commentaries (including those requiring human edits)
    \item \textbf{Implication:} Up to 25\% of accepted commentaries may require human editing; 
          up to 5\% may be rejected entirely
\end{itemize}

% =============================================================================
\section{Recommendations Priority Matrix}
\label{sec:priority-matrix}
% =============================================================================

\begin{longtable}{p{1cm} p{5cm} p{3cm} p{2cm} p{1.5cm}}
\caption{Audit Finding Resolution Status} \\
\toprule
\textbf{ID} & \textbf{Action} & \textbf{Owner} & \textbf{Timeline} & \textbf{Status} \\
\midrule
\endfirsthead
\rowcolor{green!10}
P0-1 & Reconcile state count (27→28) & Spec Owner & 1 week & \textcolor{green}{\textbf{RESOLVED}} \\
\rowcolor{green!10}
P0-2 & Add vLLM fallback specification & AI Architect & 2 weeks & \textcolor{green}{\textbf{RESOLVED}} \\
\rowcolor{green!10}
P1-1 & Define audit trail schema & Controls Owner & 4 weeks & \textcolor{green}{\textbf{RESOLVED}} \\
\rowcolor{green!10}
P1-2 & Add normality test to AC & Data Science & 4 weeks & \textcolor{green}{\textbf{RESOLVED}} \\
\rowcolor{yellow!20}
P2-1 & Automate kill-switch eval & Platform Team & 8 weeks & \textcolor{orange}{\textbf{PLANNED}} \\
\rowcolor{green!10}
P2-2 & Define commentary metrics & Product Owner & 2 weeks & \textcolor{green}{\textbf{RESOLVED}} \\
\bottomrule
\end{longtable}

\end{document}